% =====================================================================
%  IDM-JEPA architecture figure -- PAPER VERSION.
%  \input from main.tex. Do NOT add \documentclass here.
%
%  The editable source is idm_model_architecture_standalone.tex, which
%  still compiles on its own to a landscape PDF. Edit the drawing THERE,
%  then copy the tikzpicture body back into this file.
%
%  Requires in main.tex's preamble:
%    \usepackage{tikz}
%    \usetikzlibrary{positioning, arrows.meta, shapes.geometric, calc, fit}
% =====================================================================
\begin{figure}[t]
  \centering
  % ---- SIZE KNOB ----------------------------------------------------
  % The drawing is ~27cm x 10cm natural (aspect ~2.7:1), so at
  % \linewidth it is only ~5.2cm tall -- it fits a NeurIPS page without
  % crowding out the text. Lower the factor to shrink it; the height
  % follows automatically. Do not pass a height argument to \resizebox.
  \resizebox{\linewidth}{!}{%
  \begin{tikzpicture}[
      >={Latex[length=2mm, width=2mm]}, thick,
      % styles are scoped to this picture so they do not leak into the paper
      imgnode/.style={rectangle, rounded corners, draw=black!70, fill=brown!20,
                    minimum width=1.7cm, minimum height=1.7cm},
    encoder/.style={rectangle, rounded corners, draw=blue!70, fill=blue!15,
                    minimum width=2.9cm, minimum height=1.5cm, align=center,
                    font=\Large\bfseries},
    latent/.style={font=\LARGE, text=#1!60!black, inner sep=3pt},
    modelbox/.style={rectangle, rounded corners, draw=#1!80, fill=#1!15,
                    minimum width=3.4cm, minimum height=1.6cm, align=center,
                    font=\Large\bfseries},
    regnode/.style={font=\Large\bfseries, text=red!80!black, inner sep=3pt},
    lossarrow/.style={<->, dashed, draw=red!70!black, thick},
    losslabel/.style={font=\Large\bfseries, text=red!70!black},
    loop/.style={rounded corners=4pt},
  ]

  % =====================================================
  %  Layout: three horizontal bands, flowing left to right.
  %    TOP  (y= 2.6)  o_t   -> Encoder -> z_t  ------> Forward model -> zhat_{t+1}
  %    MID  (y= 0.0)  Composition -> Inverse model -> ahat_t  <~~>  a_t
  %    BOT  (y=-2.6)  o_t+1 -> Encoder -> z_t+1     Action selection -> Action encoder
  %  Natural size is ~27cm x 10cm (aspect ~2.7:1), so at \linewidth the
  %  figure is only ~5cm tall. Fonts are deliberately \Large: the picture
  %  is scaled by ~0.52, which lands them at roughly 7.5pt on the page.
  % =====================================================

  % -----------------------------------------------------
  % 1. Observations and the shared encoder (two parallel streams)
  % -----------------------------------------------------
  \node[imgnode] (ot)  at (0, 2.6) {};
  \node[left=0.15cm of ot,  font=\Large] {$o_t$};
  \node[imgnode] (ot1) at (0,-2.6) {};
  \node[left=0.15cm of ot1, font=\Large] {$o_{t+1}$};

  \node[encoder] (enc1) at (3.4, 2.6) {Encoder};
  \node[encoder] (enc2) at (3.4,-2.6) {Encoder};

  \draw[->] (ot.east)  -- (enc1.west);
  \draw[->] (ot1.east) -- (enc2.west);

  % -----------------------------------------------------
  % 2. Latents and SIGReg
  % -----------------------------------------------------
  \node[latent=blue] (zt)  at (7.0, 2.6) {$z_t$};
  \node[latent=teal] (zt1) at (7.0,-2.6) {$z_{t+1}$};

  \draw[->] (enc1.east) -- (zt.west);
  \draw[->] (enc2.east) -- (zt1.west);

  \node[regnode] (sigt)  at (4.9, 4.7) {SIGReg};
  \node[regnode] (sigt1) at (4.9,-4.7) {SIGReg};
  \draw[->] (zt.north west)  -- (sigt.south east);
  \draw[->] (zt1.south west) -- (sigt1.north east);

  % -----------------------------------------------------
  % 3. Composition function and inverse model
  % -----------------------------------------------------
  \node[modelbox=green!60!black] (comp) at (7.0,  0) {Composition\\function};
  \node[modelbox=orange]         (inv)  at (11.4, 0) {Inverse\\model};

  \draw[->] (zt.south)  -- (comp.north);
  \draw[->] (zt1.north) -- (comp.south);
  \draw[->] (comp.east) -- (inv.west);

  % -----------------------------------------------------
  % 4. Predicted vs. ground-truth action, and the selector
  % -----------------------------------------------------
  \node[latent=orange] (ahat) at (14.6, 0) {$\hat{a}_t$};
  \node[latent=black]  (at)   at (18.0, 0) {$a_t$};

  \draw[->, orange!80!black] (inv.east) -- (ahat.west);
  \draw[lossarrow] (ahat.east) -- node[above, losslabel, align=center]
                   {teacher\\forcing\\loss} (at.west);

  \node[modelbox=purple] (as)   at (11.4,-2.6) {Action\\selection};
  \node[modelbox=purple] (aenc) at (19.6,-2.6) {Action\\encoder};

  % the two candidate actions drop into the selector on parallel lanes
  \draw[->, loop, orange!80!black]
        (ahat.south) -- ++(0,-0.85) -| ([xshift=-0.8cm]as.north);
  \draw[->, loop]
        (at.south)   -- ++(0,-1.30) -| ([xshift=0.8cm]as.north);

  \draw[->] (as.east) -- (aenc.west);

  % -----------------------------------------------------
  % 5. Forward model and the latent prediction loss
  % -----------------------------------------------------
  \node[modelbox=blue] (fwd)  at (19.6, 2.6) {Forward\\model};
  \node[latent=blue]   (zhat) at (23.7, 2.6) {$\hat{z}_{t+1}$};

  \draw[->] (aenc.north) -- (fwd.south);
  \draw[->] (zt.east)    -- (fwd.west);
  \draw[->] (fwd.east)   -- (zhat.west);

  % return path for the MSE, routed below every box
  \draw[lossarrow, loop] (zhat.south) -- (23.7,-4.5)
        -- node[above, losslabel] {MSE} (7.0,-4.5) -- (zt1.south);

  \end{tikzpicture}%
  }
  \caption{% >>> Replace the bracketed sentence with the one thing you want
           % >>> the reader to take away; the rest matches the code.
           \todo{one-sentence headline}
           A shared encoder maps $o_t$ and $o_{t+1}$ to latents $z_t$ and
           $z_{t+1}$, both regularised by SIGReg. The composition function
           combines the pair and the inverse model reads it, emitting a
           predicted action $\hat a_t$ with the same width as a real
           action. Action selection takes the ground-truth $a_t$ where the
           transition is labelled and $\hat a_t$ everywhere else, so a
           single action encoder feeds the forward model in both cases and
           the downstream planner is unchanged. Training minimises the
           latent prediction error between $\hat z_{t+1}$ and $z_{t+1}$
           together with a teacher-forcing loss on $\hat a_t$, applied on
           labelled transitions only.}
  \label{fig:architecture}
\end{figure}
